\documentclass[11pt]{article}
\usepackage[margin=0.9in]{geometry}
\usepackage[T1]{fontenc}
\begin{document}
\section*{Reset to new-construction project locations}
\noindent September 11, 2026. Agent-drafted research record.

Jacob chose a simple project-location dataset: accept HUD's new-construction
classification unless repeated records or a concrete inconsistency give a reason
to investigate. The new build starts from the original HUD workbook. It uses none
of the previous physical-development links or external review decisions. The prior
workflow was removed from the active project and preserved in Git history and a
pre-reset snapshot.

\begin{center}
\small
\input{../tasks/build_lihtc/output/summary.tex}
\end{center}

The main table keeps the earliest new-construction record at each standardized
primary address. A unique record stays even when its year or hedonics are missing.
At a repeated address, a missing year or conflicting earliest-year tie prevents
selection; identical ties can use a stable representative only when their listed
characteristics agree. All raw new-construction records, including later phases,
remain in the record table. A repeated address is a review question, not proof
that a later project is a duplicate. Addresses are not parcel identifiers.

Census address-range matches provide an independent location comparison. Matches in the reported state are compared with HUD coordinates.
Non-exact strings can pass when the locations agree; without a HUD comparison,
an exact Census match is required. A discrepancy
greater than 500 meters triggers review; that screening threshold is not an error
definition. HUD-only coordinates are retained as unconfirmed. Scattered-site
projects retain their primary point without claiming that it represents every
building. Missing units or bedroom counts do not invalidate a checked location.

The remaining work is confined to repeated-address choices and unresolved
locations. There is no requirement to reconstruct rehabilitation histories or
obtain two external sources for every project. Neighborhood characteristics and
parcel matching are outside this reset. The saved Census tract codes describe
2020 geography; the analysis has not yet chosen historical neighborhood measures.

\paragraph{Reproduction.} Run \texttt{make setup}, then \texttt{make} at the
repository root. Task Makefiles are for prepared inputs; the root declares the
full dependency graph. HUD source: pinned 2024 release, retrieved August 8, 2026.
Census: Public\_AR\_ACS2025 and Census2020\_ACS2025, requests defined September 11,
2026. The code reset follows revision \texttt{fe507fd}. Saved-file reports record
rows, keys, missingness, ranges, and fingerprints. Raw geocoder responses are
retained so ordinary builds do not send the addresses again.
\end{document}
